
\documentclass{scrartcl}

% Path from this folder to the repository root, plus the two
% two worksheet-specific files (title macro and running header).
\newcommand{\npbase}{../../../}
\newcommand{\npstyle}{plain}
% ===========================================================================
%  Every \usepackage of the article documents lives here.
%
%  Loaded first by preamble.tex. style.tex and commands.tex assume everything
%  below is already available -- this file therefore has to stay self-contained
%  and free of layout decisions.
% ===========================================================================

% ---------------------------------------------------------------------------
% Build switches -- these come first because every later file keys off them.
%
%   \npstyle = screen  (default)  coloured links and info boxes, screen header
%            = print              grayscale, two-sided, print-optimised
%            = plain              no fancy header/footer at all (short handouts)
%
% A document selects one by defining \npstyle before loading the preamble.
% ---------------------------------------------------------------------------

\usepackage{etoolbox}               % \ifdefstring, \patchcmd
\usepackage{ifthen}                 % \ifthenelse, used by the macros in commands.tex

\providecommand{\npstyle}{screen}
\newif\ifnpprint \ifdefstring{\npstyle}{print}{\npprinttrue}{\npprintfalse}
\newif\ifnpfancy \ifdefstring{\npstyle}{plain}{\npfancyfalse}{\npfancytrue}

% ---------------------------------------------------------------------------
% Who wrote this. Loaded here because this is the one file every document
% pulls in, so \npauthor / \npauthorurl / \npauthormail are always available.
% ---------------------------------------------------------------------------
% ===========================================================================
%  Who wrote this and where it lives.
%
%  These three macros are the ONLY place the author's identity appears in the
%  layout (running heads, footers, title pages, slide footers).  If you fork
%  this repository, change them here and every document follows.
% ===========================================================================

\providecommand{\npauthor}{Hendrik Möller}
\providecommand{\npauthorurl}{https://github.com/Hendrik-code/numprog_script}
\providecommand{\npauthormail}{hendrik.moeller[AT]tum.de}


% ---------------------------------------------------------------------------
% Encoding, fonts and language
% ---------------------------------------------------------------------------
\usepackage[T1]{fontenc}            % proper hyphenation of words with umlauts
\usepackage[utf8]{inputenc}         % source files are UTF-8
\usepackage{lmodern}                % scalable Computer-Modern-like font
\usepackage{mathpazo}               % Palatino for the title pages
\usepackage[ngerman]{babel}         % German hyphenation and captions
\usepackage[autostyle, german=guillemets]{csquotes}
\MakeOuterQuote{"}                  % lets you type "..." for German quotes
\usepackage{relsize}

% ---------------------------------------------------------------------------
% Page layout
% ---------------------------------------------------------------------------
\usepackage{fancyhdr}               % configured in style.tex
\usepackage{multicol}
\usepackage{float}                  % [H] float placement
\usepackage{placeins}               % \FloatBarrier
\usepackage{framed}                 % the coloured info boxes
\usepackage{mdframed}
\usepackage{rotating}

% ---------------------------------------------------------------------------
% Maths
% ---------------------------------------------------------------------------
\usepackage{amsmath, accents}
\usepackage{amssymb}
\usepackage{cancel}
\usepackage{centernot}              % strike through long symbols such as ==>
\usepackage{esdiff}                 % derivative operator
\usepackage{nicefrac}
\usepackage{fp}                     % fixed point arithmetic in macros

% ---------------------------------------------------------------------------
% Tables, colour and symbols
% ---------------------------------------------------------------------------
\usepackage{booktabs}
\usepackage{colortbl}
\usepackage{xcolor}
\usepackage{enumerate}
\usepackage{pifont}                 % \ding checkmarks and crosses
\usepackage{tikzsymbols}            % \Smiley, \Cooley, \Laughey, ...
\tikzsymbolsset{global-scale=1.2}

% ---------------------------------------------------------------------------
% Graphics and plots
%
% pgfplots pulls in tikz; it MUST be loaded here because style.tex
% calls \pgfplotsset.  (Documents used to load only the style file and failed
% with "Undefined control sequence" on \pgfplotsset -- that is why there is a
% single entry point now.)
% ---------------------------------------------------------------------------
\usepackage{graphicx}
\usepackage{tikz}
\usetikzlibrary{datavisualization}
\usetikzlibrary{datavisualization.formats.functions}
\usepackage{pgfplots}
\pgfplotsset{compat=1.16}
\usetikzlibrary{pgfplots.colorbrewer}

% ---------------------------------------------------------------------------
% Code listings
% ---------------------------------------------------------------------------
\usepackage{listings}
\lstset{
    frame            = tb,
    language         = Python,
    aboveskip        = 3mm,
    belowskip        = 3mm,
    showstringspaces = false,
    columns          = flexible,
    basicstyle       = {\small\ttfamily},
    numbers          = none,
    breaklines       = true,
    breakatwhitespace= true,
    tabsize          = 3,
}

% ---------------------------------------------------------------------------
% Programming helpers used by commands.tex
% ---------------------------------------------------------------------------
\usepackage{environ}                % \NewEnviron
\usepackage{catchfilebetweentags}   % pulling snippets into the complete script
\usepackage{nccfoots}               % \Footnotetext
\usepackage{verbatim}
\usepackage{hyperref}               % keep last: it patches other packages

% ===========================================================================
%  Macros for the Trainingsblaetter.
%
%  Self-contained on purpose: the worksheets have their own title page
%  (\HAtitle) and their own task/solution layout, and share only the package
%  list with the rest of the repository (gemeinsam/preamble/packages.tex).
% ===========================================================================

\newcommand{\HAtitle}[2]{
	\begin{titlepage} % Suppresses displaying the page number on the title page and the subsequent page counts as page 1
		\newcommand{\HRule}{\rule{\linewidth}{0.5mm}} % Defines a new command for horizontal lines, change thickness here
		\center % Centre everything on the page
		\hspace{1cm} \\[3.5cm]
		\HRule\\[0.4cm]
		{\huge\bfseries Numerisches Programmieren Trainingsaufgaben #1}\\[0.4cm] % Title of your document
		\textsc{\Large #2}\\[0.5cm] % Minor heading such as course title
		\HRule\\[1.5cm]
		\begin{minipage}{0.4\textwidth}
			\begin{flushleft}
				\large
				\textit{Autor}\\
				\textsc{\npauthor{}} % Your name
			\end{flushleft}
		\end{minipage}
		~
		\begin{minipage}{0.4\textwidth}
			\begin{flushright}
				\large
				\textit{Version vom:}\\
				\textsc{\today} % Supervisor's name
			\end{flushright}
		\end{minipage}
		\vfill\vfill % Position the date 3/4 down the remaining page
		\large{\textit{Erst die Aufgaben rechnen, bevor Ihr die Lösung anguckt!}}
		\vfill % Push the date up 1/4 of the remaining page
	\end{titlepage}
}

\definecolor{easy}{rgb}{0.2 0.2 1}
\definecolor{medium}{rgb}{0.9 0.4 0}
\definecolor{hard}{rgb}{1 0 0.4}
\definecolor{special}{rgb}{0.5 0.2 1}
\newcommand{\easy}{\textcolor{easy}{(Easy)}}
\newcommand{\medium}{\textcolor{medium}{(Medium)}}
\newcommand{\hard}{\textcolor{hard}{(Hard)}}
\newcommand{\transfer}{\textcolor{special}{(Transfer)}}

\newcommand{\itemf}{\item}
\newcommand{\itemeasy}{\item[\color{easy}]\alph*)}


\newcommand{\task}[4]{
	\clearpage
	\section{#1}
	%
	#2
	\begin{enumerate}[label=\alph*)]
		#3
	\end{enumerate}
	
	\vfill \textit{Auf der nächsten Seite folgt die Lösung\dots}
	\clearpage
	\Large \textbf{Lösung:}
	\vspace{0.3cm}
	\normalsize
	\begin{enumerate}[label=\alph*), nolistsep, itemsep=\fill]
		#4
	\end{enumerate}
	\vfill
}

\newcommand{\subtask}[1]{
	\begin{enumerate}[label=(\roman*)]
		#1
	\end{enumerate}
}

\newcommand{\hinweis}{\textit{Hinweis: }}
% =========================================================
% Specific math variables:
% =========================================================

\newcommand{\fabs}{f_{\text{abs}}}
\newcommand{\frel}{f_{\text{rel}}}

\newcommand{\maxError}{\epsilon_{\text{Ma}}}

\newcommand{\mddots}{\reflectbox{$\ddots$}}


\newcommand{\umax}{_\text{max}}
\newcommand{\umin}{_\text{min}}

% =========================================================
% Math Commands
% =========================================================

% Should be renamed or removed
\newcommand{\st}[1]{^{(#1)}}

% vertical line in the middle:
\newcommand{\separator}[1]{\; #1 \;}
\newcommand{\sep}{\separator{|}}
\newcommand{\sepArrow}{\separator{\Longleftrightarrow}}

% BigO
\newcommand{\bigO}{\mathcal{O}} % Big-O notation

% Norm with double |

% Matrix Transpose and inverted
\newcommand{\trans}{^{\mathrm{\mathsmaller T}}}
\newcommand{\inv}{^{\mathrm{\mathsmaller{\text{-}1}}}}
\newcommand{\unitM}{\mathbb{1}}


% informal descriptions:
\newcommand{\uleftDown}{_\text{linksUnten}}
\newcommand{\uleftUp}{_\text{linksOben}}
\newcommand{\uleft}{_\text{links}}
\newcommand{\uleftDiagonal}{_\text{ganzLinksUnten}}
\newcommand{\uleftUpDiagonal}{_\text{ganzLinksOben}}

% Soll sein
\newcommand{\shouldEq}{ \stackrel{!}{=} }

% Formelnamen:
\newcommand{\QT}{Q_{\text{T}}}
\newcommand{\QTS}{Q_{\text{TS}}}
\newcommand{\QS}{Q_{\text{S}}}
\newcommand{\QSS}{Q_{\text{SS}}}
\newcommand{\RTS}{R_{\text{TS}}}
\newcommand{\RSS}{R_{\text{SS}}}
%--------------------------

% correct strange spacing around delimiters
\renewcommand{\l}{\mathopen{}\mathclose\bgroup\left}
\renewcommand{\r}{\aftergroup\egroup\right}

% norms
\newcommand{\norm}[1]{\l\lVert #1 \r\rVert}
\newcommand{\pnorm}[2]{\norm{#1}_{#2}}

\newcommand{\nvector}[1]{
	\begin{pmatrix}
		#1
	\end{pmatrix}
}

\newcommand{\nmatrix}[1]{
	\begin{bmatrix}
		#1
	\end{bmatrix}
}

\def\nLongrightarrow{\centernot \Longrightarrow}
\def\nLongleftrightarrow{\centernot \Longleftrightarrow}
\def\nforall{\centernot \forall}
\def\nexists{\centernot \exists}
\def\nsqsubseteq{\centernot \sqsubseteq}
\def\nsubset{\centernot \subset}
\def\nequiv{\centernot \equiv}

\def\symdif{\triangle}
\def\imp{\rightarrow}
\def\bik{\leftrightarrow}
\def\xor{\oplus}
\def\nand{\, \bar{\land}\,}
\def\nor{\,\bar{\lor}\,}


% number sets
\newcommand{\booleans}[0]{\mathbb{B}}			% Perhaps call it \binaries and defined it as \residuals{2} instead.
\newcommand{\naturals}[0]{\mathbb{N}}
\newcommand{\wholes}[0]{\mathbb{N}_0}
\newcommand{\integers}[0]{\mathbb{Z}}
\newcommand{\rationals}[0]{\mathbb{Q}}
\newcommand{\reals}[0]{\mathbb{R}}
\newcommand{\complexes}[0]{\mathbb{C}}
\newcommand{\residuals}[1]{\mathbb{Z}_{#1}}
\newcommand{\primes}[0]{\mathbb{P}}

% brackets (round, square, curly, angled)
\renewcommand{\(}[0]{\l(}						% \(...\) was defined for $...$ as \[...\] for $$...$$
\renewcommand{\)}[0]{\r)}
\newcommand{\rb}[1]{\l(#1\r)}
\renewcommand{\sb}[1]{\l[#1\r]}
\newcommand{\cb}[1]{\l\{#1\r\}}
\newcommand{\ab}[1]{\l\langle#1\r\rangle}

% generating set
\newcommand{\generatingset}[1]{\ab{#1}}

% floor and ceiling
\newcommand{\ceil}[1]{\l\lceil#1\r\rceil}
\newcommand{\floor}[1]{\l\lfloor#1\r\rfloor}


% absolute value
\newcommand{\abs}[1]{\l|#1\r|}

% integrals
\newcommand{\integral}[4]{\int_{#1}^{#2} #3 \ \mathrm{d}#4}
\newcommand{\antiderivative}[2]{\integral{}{}{#1}{#2}}
\renewcommand{\d}[1]{\ \mathrm{d}#1}

% logarithms
\newcommand{\ld}[0]{\operatorname{ld}}
\newcommand{\lb}[0]{\operatorname{lb}}

% signum function
\newcommand{\sgn}[0]{\operatorname{sgn}}

% fixed set
\newcommand{\fix}[0]{\operatorname{fix}}		% fix(f) = {x|f(x) = x}

% closures
\newcommand{\refl}[0]{\operatorname{ref}}
\newcommand{\symm}[0]{\operatorname{sym}}

% sets of functions
\newcommand{\inj}[0]{\operatorname{inj}}
\newcommand{\surj}[0]{\operatorname{surj}}
\newcommand{\bij}[0]{\operatorname{bij}}

% fields
\newcommand{\charac}[0]{\operatorname{char}}
\newcommand{\GF}[0]{\operatorname{GF}}

% least common multiple and greatest common divisor
\newcommand{\ggT}[0]{\operatorname{ggT}}
\newcommand{\kgV}[0]{\operatorname{kgV}}

% average
\newcommand{\avg}[0]{\operatorname{avg}}

% order
\newcommand{\ord}[0]{\operatorname{ord}}

% ----------
% round 
\newcommand{\rd}[0]{\operatorname{rd}}
\newcommand{\rdM}[0]{\operatorname{rd_M}}

% condition value
\newcommand{\cond}[0]{\operatorname{cond}}

% complex numbers
\newcommand{\realPart}{\operatorname{Re}}
\newcommand{\imaginaryPart}{\operatorname{Im}}

% Fourier transformation
\newcommand{\DFT}{\operatorname{DFT}}
\newcommand{\IDFT}{\operatorname{IDFT}}
\newcommand{\FFT}{\operatorname{FFT}}
\newcommand{\BFO}{\operatorname{BFO}}

% diagonal entries of a matrix
\newcommand{\diag}{\operatorname{diag}}
% --------

% sequent calculus
\newcommand{\FL}[0]{\operatorname{FL}}
\newcommand{\Nat}[0]{\operatorname{Nat}}

% true and false
\newcommand{\false}[0]{\operatorname{false}}
\newcommand{\true}[0]{\operatorname{true}}

% expected value
\newcommand{\Eempty}[0]{\mathbb{E}}
\newcommand{\E}[1]{\Eempty\sb{#1}}
\newcommand{\Eindex}[2]{\Eempty_{#1}\sb{#2}}
\newcommand{\Econd}[2]{\Eempty\sb{#1 \middle| #2}}	% E[#1|#2]

% variance
\newcommand{\Varempty}[0]{\operatorname{Var}}
\newcommand{\Var}[1]{\Varempty\sb{#1}}
\newcommand{\Varindex}[2]{\Varempty_{#1}\sb{#2}}
\newcommand{\Varcond}[2]{\Varempty\sb{#1 \middle| #2}}	% Var[#1|#2]

\newcommand{\MSE}[0]{\operatorname{MSE}}
\newcommand{\Bias}[0]{\operatorname{Bias}}

% probability distributions
\newcommand{\Ber}[0]{\operatorname{Ber}}
\newcommand{\Bin}[0]{\operatorname{Bin}}
\newcommand{\Poi}[0]{\operatorname{Poi}}
\newcommand{\Geo}[0]{\operatorname{Geo}}
\newcommand{\NegBin}[0]{\operatorname{NegBin}}
\newcommand{\Hyp}[0]{\operatorname{Hyp}}
\newcommand{\NegHyp}[0]{\operatorname{NegHyp}}
\newcommand{\Nor}[0]{\operatorname{\mathcal{N}}}
\newcommand{\Exp}[0]{\operatorname{Exp}}
\newcommand{\Erl}[0]{\operatorname{Erl}}
\newcommand{\ConUni}[0]{\operatorname{Uni}}
\newcommand{\DisUni}[0]{\operatorname{Uni}}
\newcommand{\Uni}[0]{\operatorname{Uni}}					% TODO: cambiar por \UniDis y \UniCon

% relations
\newcommand{\id}[1]{\operatorname{id}_{#1}}					% Identitätsrelation
\newcommand{\all}[1]{#1 \times #1}							% Allrelation

% =========================================================
% TikZ Commands
% =========================================================

\newcommand{\tikzPoint}[3]{
	\node[label={#1:{#2}},circle,fill,inner sep=2pt] at (axis cs:#3) {};
}
\newcommand{\tikzPointLeft}[2]{
	\tikzPoint{180}{#1}{#2}
}
\newcommand{\tikzPointRight}[2]{
	\tikzPoint{0}{#1}{#2}
}
\newcommand{\tikzPointUp}[2]{
	\tikzPoint{90}{#1}{#2}
}
\newcommand{\tikzPointDown}[2]{
	\tikzPoint{270}{#1}{#2}
}
\newcommand{\tikzPointUL}[2]{
	\tikzPoint{135}{#1}{#2}
}
\newcommand{\tikzPointUR}[2]{
	\tikzPoint{45}{#1}{#2}
}
\newcommand{\tikzPointDL}[2]{
	\tikzPoint{225}{#1}{#2}
}
\newcommand{\tikzPointDR}[2]{
	\tikzPoint{315}{#1}{#2}
}

%[red, very thick, |-] (0,0) -- (22,0);
\def \tikzPattern{on 1pt off 3pt on 3pt off 3pt}%dash pattern=on 1pt off 3pt on 3pt off 3pt}

\newcommand{\tikzLineDotted}[3]{
	\draw[gray, dotted] (axis cs:#2) -- node[left]{#1} (axis cs:#3);
}
\newcommand{\tikzLine}[4]{
	\draw[#1] (axis cs:#3) -- node[left]{#2} (axis cs:#4);
}
\newcommand{\tikzLineUp}[4]{
	\draw[#1] (axis cs:#3) -- node[above]{#2} (axis cs:#4);
}
\newcommand{\tikzNormalLine}[3]{
	\tikzLine{black, thick}{#1}{#2}{#3}
}
\newcommand{\tikzFatLine}[3]{
	\tikzLine{blue, very thick}{#1}{#2}{#3}
}
\newcommand{\tikzFatterLine}[3]{
	\tikzLine{blue, line width=0.65mm}{#1}{#2}{#3}
}
\newcommand{\tikzMegaFatLine}[3]{
	\tikzLine{blue, line width=0.95mm}{#1}{#2}{#3}
}

\newcommand{\tikzArrow}[3]{
	\draw[very thick, ->] (axis cs:#2) -- node[above]{#1} (axis cs:#3);
}
% Angles:
% 0 -> right, 90 = up, 270 = down, 180
\newcommand{\tikzCircle}[3]{
	\draw(axis cs:#2) circle[#1, radius=#3];
}

\newcommand{\tikzTrigonometric}{
	xtick={-3.14159, -1.5708, 0,1.5708,3.14159,4.7123889},
	xticklabels={$-\pi$, $-\frac{\pi}{2}$, $0$,$\frac{\pi}{2}$,$\pi$,$\frac{3\pi}{2}$},
}


\newcommand{\tikzFigure}[2]{
	\begin{figure}[ht]
		\centering
		\begin{tikzpicture}
		#1
		\end{tikzpicture}
		\caption{#2}
	\end{figure}
}
\newcommand{\tikzFigureStandard}[4]{
	\tikzFigure{
		\begin{axis}[
			axis lines=left,
			xlabel=$x$,
			ylabel=$y$,
			enlargelimits,
			%xtick={-3.14159, -1.5708, 0,1.5708,3.14159,4.7123889},
			%xticklabels={$-\pi$, $-\frac{\pi}{2}$, $0$,$\frac{\pi}{2}$,$\pi$,$\frac{3\pi}{2}$},
			%ymax=2,
			%ymin=-0.5]
			%[domain=0:1,legend pos=outer north east]
			\addplot[domain=#1:#2] {#3} node[left]{};
			%				( 0, 4 / pi )
			%				( pi, 0 )
			%				( 0, 0)
			%				( 0, 4 / pi)
			%			};
		\end{axis}
	}{#4}
}
% ===========================================================================
%  Page style for the Trainingsblaetter.
%
%  The worksheets load gemeinsam/preamble/packages.tex from the root for all
%  packages, and then this file for their own running header: it shows the
%  sheet number (\num) and the topic (\topic) that each worksheet defines.
% ===========================================================================

\usepackage{fancyhdr}
\pagestyle{fancy}
\fancyhf{}
\fancyhead[C]{\npauthor{} -- \num \;-- \topic}
\fancyfoot[C]{\textbf{\thepage}}
\renewcommand{\headrulewidth}{1pt}
\renewcommand{\footrulewidth}{1pt}

% Worksheets use blank-line paragraph spacing instead of indentation, and
% enumitem for the (a) (b) (c) sub-task labels.
\usepackage[skip=0.0cm, indent=0pt]{parskip}
\usepackage{enumitem}

\usepackage{hyperref}
\hypersetup{
	colorlinks = true,
	linktoc    = all,
	linkcolor  = blue,
	urlcolor   = blue,
}



\usepackage[left=3.1cm, right=3.1cm, bottom=3cm, top=3cm]{geometry}
% =========================================================
\newcommand{\num}{01}
\newcommand{\topic}{Zahlendarstellung}


\begin{document}
	\HAtitle{\num}{\topic}
	\newpage
%--------------------------------------------------
	\task{IEEE Umwandlung Schritt für Schritt}
	{
		Wir haben in den Tutorien das IEEE Zahlenformat kennengelernt. Es nutzt $32$ Bits: Eins für ein Vorzeichen, $8$ für den Exponenten und die restlichen $23$ für die sogenannte Mantisse. Die folgende Aufgabe soll die schrittweise Umwandlung einer dezimalen Zahl zum IEEE Format vertiefen.
	}{
		\itemf Wandeln Sie $14$ in das binäre System um.
		\itemf Überlegen sie aus (a), welcher (dezimale) Exponent benötigt wird, damit eine Mantisse normiert ist.
		\itemf Leiten Sie nun nacheinander aus (b) erst das Vorzeichenbit, dann die Exponentenbits und dann die Mantisse her.
		\itemf Wandeln Sie nun $-\frac{1}{6}$ in das IEEE Format um.
	}{
		\itemf $14$ lässt sich binär herunterbrechen auf $8 + 4 + 2 = 14$. Somit ist die binäre Zahl $1110$.
		\itemf Die binäre Zahl ist $1110$. Eine Mantisse ist genau dann normiert, wenn diese eine führende Eins hat und sofort ein Komma folgt. Wir können also mit folgender Rechnung herleiten:
		\begin{align*}
			1110 = 111.0 \cdot 2 = 111.0 \cdot 2^1 = 11.10 \cdot 2^2 = 1.110 \cdot 2^3
		\end{align*}
		Der Exponent ist also eine Kommaverschiebung nach links. Ein negativer Exponent wäre entsprechend eine Verschiebung nach rechts. Der benötigte Exponent für die Normierung ist hier also $3$, da hierbei die Mantisse $1.110$ entspricht. Eine führende Eins und direkt danach ein Komma.
		\itemf Vorzeichenbit ist das einfachste: Die Zahl $14$ ist positiv, also ist das Vorzeichenbit $0$. 
		
		Der dezimale Exponent ist $3$, wie wir aus (b) wissen. Auf diesen müssen wir den IEEE Offset von 127 addieren. Wir erhalten $130$. Diese Zahl müssen wir jetzt nur noch in Binär darstellen. Das ist entsprechend $10000010$.
		
		Die Mantisse können wir direkt aus (b) auslesen, sie ist $1.110$.
		
		Wenn wir jetzt alles zusammensetzen, müssen wir nur aufpassen, dass wir von der Mantisse nicht die führende Eins oder das Komma mitnehmen, da diese bei IEEE nicht abgespeichert wird. Wir erhalten (mit Reihenfolge Vorzeichen|Exponent|Mantisse):
		\begin{align*}
			0|100000010|11000000000000000000000
		\end{align*}
		\itemf Erstmal in binär umwandeln. Das ist nicht trivial, also muss man schriftliche Division im binären Bereich anwenden.
		$-\frac{1}{6} = -001_2 : 110_2$
		\begin{table}[H]
			\centering
			\label{table:dec>binary1}
			\begin{tabular}{lllllll}
				1 & & & & & : $110_2$ &= \\
				1 & 0 & 0 & 0 & & &$=0.00$\\
				- & 1 & 1 & 0 & & &$=0.001$\\ \cline{2-5}
				  & 0 & 1 & 0 & 0 & &\\
				  &  & 1 & 0 & 0 & 0 &$=0.0010$\\
				  & - & 0 & 1 & 1 & 0 &$=0.00101$\\ \cline{3-6}
				  &   &   &   & 1 & 0 &\\
				  & & & & & \dots 
			\end{tabular}
		\end{table}
		$-\frac{1}{6}$ ist also $0.0\overline{01}_2$. Der benötigte Exponent ist $-3$, da dann die Mantisse $1.\overline{01}_2$ und damit normiert ist. Jetzt noch in IEEE Format umwandeln.
		
		Vorzeichenbit ist $1$, da die Zahl negativ ist.
		
		Exponent ist $-3+127 = 124$ in binär, also $01111100$. 
		
		Die Mantisse können wir ablesen, die Periode müssen wir solange wiederholen, bis uns die Bits ausgehen. Also $01010101010101010101010|101$. Wir müssen aber alles nach den 23 Bits (hier mit einem $|$ markiert) runden. Hier haben wir den eindeutigen Fall, bei dem aufgerundet wird.
		
		Vollständig haben wir also: $1|01111100|01010101010101010101011$.
	}
%--------------------------------------------------
	\task{Gleitkomma Arithmetik}
	{
		Nun gucken wir uns das Thema der Maschinenzahlen etwas allgemeiner an.
	}{
		\itemf Was ist der Nachteil eines Vorzeichenbits? Warum entsteht dieser Nachteil?
		\itemf Nehmen Sie eine binäre Darstellung mit 3 signifikanten Stellen an. Formulieren Sie eine Rechnung zwischen 2 Zahlen ihrer Wahl, bei dem das Problem der Absorption auftreten wird. Berechnen Sie von Ihrer Rechnung den absoluten sowie relativen Fehler.
		\itemf Berechnen Sie $101.11 - 110.011$ bei einer binären Darstellung mit 3 signifikanten Stellen. Runden Sie immer wie in der Übung besprochen. Was passiert hier? Berechnen Sie den relativen Fehler dieser Rechnung.
		%Sei $x=10.01_2$ und $y=111_2$ und wir haben 3 signifikante Stellen zur Verfügung. Berechnen Sie das Ergebnis von $x+y$ und bestimmen Sie den absoluten sowie relativen Fehler.
	}{
		\itemf Das Vorzeichenbit bringt meistens den Nachteil der uneindeutigen Null. Es gibt eine bestimmte Kombination aus Exponent und Mantisse (meist nur Nullen), die für den Wert der $0$ steht. Dadurch kann man sowohl eine $+0$ als auch eine $-0$ darstellen, obwohl diese vom Wert her identisch sind.
		\itemf Hier kann man viele Beispiele bringen. Am einfachsten ist es zwei Zahlen zu wählen, die mehr als die drei Stellen auseinander liegen. Also ein Beispiel:
		\begin{align*}
			100 + 0.011 = 100.011
		\end{align*}
		Hier würde $100$ die $0.011$ absorbieren, das eigentlich exakte Ergebnis $100.011$ ($4.375$ in dezimal) muss gerundet werden, da es mehr als 3 signifikante Stellen hat. Hier wird abgerundet, da die erste nicht darstellbare Ziffer eine $0$ ist. Das gerundete Ergebnis ist $100$. Damit ist die Addition also völlig wertlos gewesen!
		
		Der absolute Fehler berechnet sich durch $\abs{x - \rd(x)}$. In dem Fall also $100.011 - 100 = 0.011$. Im dezimalen Bereich entspräche das also $\frac{3}{8}$.
		
		Um den relativen Fehler zu erhalten müssen wir den absoluten nochmal durch das exakte Ergebnis teilen, in dem Fall also $\frac{\nicefrac{3}{8}}{4.375} = 0.0857$, also etwa $8.6\%$ Abweichung.
		\itemf Vor der Berechnung müssen wir beide Zahlen runden, da sie mehr als 3 signifikante Stellen haben. $101.11$ rundet sich auf zu $110$ und $110.011$ rundet man ab auf $110$. Das Ergebnis der Subtraktion ist trivialerweise $0$. Der exakte Wert wäre aber $101.11 - 110.011 = -0.011$ ($-0.375$ in dezimal). Der relative Fehler ist logischerweise immens!
	}
%--------------------------------------------------
	\task{Fremde Zahlenformate}
	{
		In dieser Aufgabe werden fiktive Zahlendarstellungen vorgestellt und dazu Verständnisfragen gestellt. Es handelt sich also um Transferaufgaben.
	}{
		\itemf Ein Zahlenformat, hier mal mit NDR abgekürzt, verwendet ein Vorzeichenbit, 3 Bits für einen Exponenten, wobei dieser vorzeichenbehaftet ist (das erste dieser 3 Bits steht also für das Vorzeichen des Exponentens) und sonst einer Ganzzahl entspricht. Für die Mantisse stehen 4 Bits zur Verfügung. Diese beginnt immer mit einer führenden Eins und ist wie bei IEEE normiert. Allerdings muss die führende Eins bei NDR abgespeichert werden.
		
		\hinweis $0|101|1000$ im NDR entspräche $2^{-1} \cdot 1.000 = 0.1 = \frac{1}{2}$
		\subtask{
			\itemf Was ist die größte und kleinste darstellbare Zahl mit dem NDR Format?
			\itemf Bestimmen Sie die Maschinengenauigkeit von NDR.
			\itemf Berechnen Sie die betragsmäßig kleinste Zahl $<0$ des NDR Formates.
		}
		\itemf Das CNF (crazy numbering format) ist folgendermaßen aufgebaut: Die Mantisse besteht aus 3 Bits. Allerdings muss diese nicht normiert sein, sondern wird direkt als binäre Zahl interpretiert. Der Exponent besteht auch aus 3 Bits und wird im 2-Komplement Format dargestellt. Es gibt kein Vorzeichenbit.
		
		\hinweis $001|011$ entspräche $2^1 \cdot 011 = 110 = 6$.
		\subtask{
			\itemf Ist das CNF Format eindeutig? Machen Sie ein Beispiel, dass Ihre These bekräftigt.
			\itemf Geben Sie die kleinste Zahl $>0$ an, die mit dem CNF gerade nicht mehr darstellbar ist.
			\itemf Wieviele verschiedene Zahlen kann das CNF wirklich repräsentieren? \hard
		}
	}{
		\itemf 
		\subtask{
			\itemf Die kleinste darstellbare Zahl erhalten wir, wenn wir das Vorzeichen auf $1$ setzen, die Zahl also negativ wird und ansonsten probieren die betragsmäßig größte Zahl zu erhalten. Das wäre bei $1|011|1111$ der Fall. Wichtig hier ist, dass das erste Bit des Exponenten keine Eins ist, sonst wäre der Exponent negativ und würde unsere Zahl betragsmäßig klein machen. Obige Darstellung entspräche übersetzt $-1 \cdot 2^3 \cdot 1.111 = -1111$. Dezimal wäre das eine $-15$.
			
			Die größte darstellbare Zahl machen wir analog. $0|011|1111$ ist $2^3 \cdot 1.111 = 1111$ und damit eine $+15$.
			\itemf Die Maschinengenauigkeit beschreibt die größte positive Zahl, die mit Eins addiert gerundet wieder Eins ergibt (also genau weg absorbiert wird). Also müssen wir uns überlegen, wie klein diese Zahl sein muss. Das ist bei NDR genau dann der Fall, wenn unsere Mantisse nicht mehr ausreicht, um die Zahl darzustellen. Schon mit der addierten Eins eingerechnet muss unsere Mantisse also $1.0001$ entsprechen, die gesuchte Zahl ergo $0.0001$. Das ist äquivalent zu $1 \cdot 2^{-4} = \frac{1}{16}$.
			
			Wir können auch überprüfen, dass das stimmt, in dem wir den Weg rückwärts gehen. Wir legen die Maschinengenauigkeit als $0.0001$ fest. Also muss $0.0001 + 1.0$ gerundet im NDR Format Eins ergeben. Wir probieren $1.0001$ in NDR darzustellen. Das wäre $0|000|1000 1$. Die Mantisse muss gerundet werden, und hier haben wir den Fall, bei dem abgerundet wird. Also erhalten wir $0|000|1000$, das im dezimalen der Eins entspricht. Die Maschinengenauigkeit kann auch nicht größer sein, da sobald wir unsere Zahl etwas erhöhen, z.B. $000100000001$, diese bei der Rechnung aufgerundet wird und somit eine Zahl größer Eins herauskommen würde.
			\itemf Wir suchen eine Zahl $<0$, also ist das Vorzeichenbit eine Eins, die gesuchte Zahl ist negativ.
			
			Betragsmäßig kleinste Zahl, also muss der Exponent so gering wie möglich sein. Das ist bei $111$ der Fall, da das einer $-3$ entspricht.
			
			Die Mantisse muss eine führende Eins haben, aber weitere Einsen würden die Zahl betragsmäßig nur größer machen, also setzen wir sie auf $1000$.
			
			Insgesamt haben wir also $1|111|1000$. Das entspricht $-1 \cdot 2^{-3} \cdot 1.000 = -0.001$ und das wiederum der $-\frac{1}{8}$.
		}
		\itemf  
		\subtask{
			\itemf Nein, das CNF ist keinesfalls eindeutig. Man betrachte die Zahlen $001|010$ und $010|001$ an. Übersetzt erhält man $2^1 \cdot 10 = 100 = 2^2 \cdot 001 = 010|001$. Beide repräsentieren die Zahl $4$.
			\itemf Fangen wir mit der kleinsten Zahl $>0$ an, die darstellbar ist mit CNF. Das wäre $100|001$. Das entspricht nämlich $2^{-4} \cdot 1 = 2^{-4}$. Jetzt suchen wir nur noch die nächstkleinere Zahl, die eben nicht mehr darstellbar ist. Der Twist ist, dass wir jetzt nicht als Mantisse $0001$ wählen können, da diese ja direkt als binäre Zahl interpretiert wird. Nur das Verringern des Exponenten kann die Zahl noch kleiner machen. Die nächstkleinere Zahl, die also nicht mehr darstellbar wäre, ist $2^{-5}$.
			\itemf Rein theoretisch würde man die Anzahl der verschiedenen Bit-Kombinationen zählen. In diesem Fall haben wir 6 Bits, es gibt also $2^6 = 64$ verschiedene Kombinationen, wie die Bits gesetzt werden können. Leider stimmt das aber nicht, da wir durch die fehlende Normierung sehr viele Kombinationen haben, die zu demselben Wert führen, siehe dazu (i). Wir können auch nicht einfach die kleinste und größte Zahl anschauen und davon auf die Anzahl Zahlen schließen, da es einige Lücken gibt. Zum Beispiel können wir die Zahl $9$ nicht darstellen, jedoch aber die $8$ und die $10$.
			Es gibt aber immer nur 8 Möglichkeiten für die Mantisse und diese repräsentieren die Werte von $0$ bis $7$ in ganzen Schritten. Das in Kombination mit den verschiedenen Exponenten hilft uns weiter.
			
			Nehmen wir einen Exponent von $000$ an, also $2^0$. Das gibt uns also die Zahlen von $0$ bis $7$ (mit den verschiedenen Kombinationen der Mantisse). Der Exponent $001$ gibt mit den verschiedenen Mantissen die Zahlen $\{0, 2, 4, 6, 8, 10, 12, 14\}$. Wenn wir alle durchgehen:
			\begin{table}[H]
				\centering
				\label{table:numbers}
				\begin{tabular}{ccll}
					Exponent & Faktor & Mögliche Zahlen & Neu hinzugekommen\\ \toprule
					$000$ & 1 & $\{0, 1, 2, 3, 4, 5, 6, 7\}$ & $\{0, 1, 2, 3, 4, 5, 6, 7\}$\\
					$001$ & 2 & $\{0, 2, 4, 6, 8, 10, 12, 14\}$ & $\{8, 10, 12, 14\}$\\
					$010$ & 4 & $\{0, 4, 8, 12, 16, 20, 24, 28\}$ & $\{16, 20, 24, 28\}$\\
					$011$ & 8 & $\{0, 8, 16, 24, 32, 40, 48, 56\}$ & $\{32, 40, 48, 56\}$
				\end{tabular}
			\end{table}
			Spätestens hier kann man jetzt schon ein Muster erkennen. Jeder Exponent fügt 4 neue mögliche Zahlen hinzu, da alle anderen schon durch den vorherigen Exponenten darstellbar sind. Das können wir auch mit der Mantisse aus 3 Bits erklären. Aus den 8 Möglichkeiten der Mantisse gibt es 4, die man redundant mit dem Exponenten ändern kann. Das wäre einmal $000$, da Null mal einen Faktor immer noch Null ergibt. Die anderen sind $001$, die man durch einen Exponentenänderung äquivalent zu der Mantisse $010$ machen kann, wie in (i). $010$ kann man zu $100$ shiften und $011$ zu $110$.
			
			So wird das also für die negativen Exponenten genauso laufen. Wir können also jetzt schon berechnen, wieviele verschiedene Zahlen das CNF darstellen kann. Zur Verständlichkeit halber liste ich hier trotzdem mal ein Beispiel auf:
			\begin{table}[H]
				\centering
				\label{table:numbers2}
				\begin{tabular}{ccll}
					Exponent & Faktor & Mögliche Zahlen & Neu hinzugekommen\\ \toprule
					$111$ & $\nicefrac{1}{2}$ & $\{0, \frac{1}{2}, 1, \frac{3}{2}, 2, \frac{5}{2}, 3, \frac{7}{2}\}$ & $\{\frac{1}{2}, \frac{3}{2}, \frac{5}{2}, \frac{7}{2}\}$\\
					$110$ & $\nicefrac{1}{4}$ & $\{0, \frac{1}{4}, \frac{1}{2}, \frac{3}{4}, 1, \frac{5}{4}, \frac{3}{2}, \frac{7}{4}\}$ & \dots \\
					$101$ & $\nicefrac{1}{8}$ & \dots & \\
					$100$ & $\nicefrac{1}{16}$ & &
				\end{tabular}
			\end{table}
			Insgesamt haben wir also die Start 8 Zahlen (von 0 bis 7) plus sieben verschiedene mögliche Exponenten, die alle genau 4 neue Zahlen hinzufügen.
			
			$8 + 7 \cdot 4 = 36$. CNF kann genau $36$ unterschiedliche Zahlen darstellen, obwohl es  $64$ unterschiedliche Kombinationen der Bits existieren. Man kann also auf $64-36 = 28$ redundante Zahlendarstellungen schließen!
		}
	}
\end{document}